\chapter{Thiết kế Hệ thống và Triển khai Kỹ thuật}
\label{ch:system_design_implementation}

Chương này trình bày chi tiết thiết kế kỹ thuật, mô hình hóa toán học và quy trình triển khai phần mềm của khung bảo mật nhận thức hai tầng SENTINEL. Được cấu trúc như một quy trình phát hiện đe dọa và phản hồi sự cố lai (hybrid), hệ thống này được thiết kế nhằm san lấp khoảng cách hiệu năng giữa xử lý gói tin mạng ở tốc độ đường truyền (wire-speed) và khả năng suy luận tiêu tốn tài nguyên của Mô hình Ngôn ngữ Lớn (LLM), đồng thời thiết lập các rào chắn nhận thức (cognitive guardrails) mạnh mẽ để cách ly LLM khỏi các hành vi thao túng đối kháng. Nội dung chương tập trung phân tích sơ đồ khối hệ thống tổng quát, các thuật toán tiền lọc thống kê trực tuyến, quy trình điều phối nhận thức dựa trên đồ thị trạng thái, cơ chế truy xuất kép Dual-RAG, lược đồ cơ sở dữ liệu Bộ nhớ Đe dọa SQLite, rào chắn bảo mật AI, chuỗi liên kết log bảo mật HMAC, giao diện người dùng dashboard Streamlit và mô hình triển khai ảo hóa Docker.

\section{Tổng quan về Kiến trúc Nhận thức Hai tầng}

Để bảo vệ môi trường mạng doanh nghiệp mà không phải đánh đổi bằng độ trễ tính toán quá lớn, hệ thống SENTINEL chia quy trình phát hiện và phản hồi mối đe dọa thành hai tầng hoạt động chuyên biệt và riêng biệt. Sự phân tách trách nhiệm kiến trúc (architectural separation of concerns) này được minh họa trực quan thông qua sơ đồ khối hệ thống tổng thể trong Hình~\ref{fig:soc_ai_concept}.

\begin{figure}[!htbp]
    \centering
    \includegraphics[width=0.95\textwidth,height=0.72\textheight,keepaspectratio]{images/soc_ai_concept.png}
    \caption{Sơ đồ kiến trúc tổng thể của hệ thống SENTINEL, minh họa luồng dữ liệu hai tầng và các thành phần cấu trúc. (Nguồn: Tác giả)}
    \label{fig:soc_ai_concept}
\end{figure}

Kiến trúc đề xuất phân chia trách nhiệm trên hai tầng chính. Tầng 1, lớp lọc dữ liệu phi trạng thái ở tốc độ đường truyền, được bố trí trực tiếp tại ranh giới mạng; tại đây, nó nạp dữ liệu đo lường tốc độ cao như nhật ký luồng và siêu dữ liệu gói tin, rồi hoặc tính toán độ lệch thống kê thời gian thực bằng thuật toán Welford, hoặc áp dụng chữ ký regex để loại bỏ hay ghi nhật ký các sự kiện lành tính ngay lập tức ở tốc độ đường truyền. Tầng 2, lớp suy luận nhận thức có trạng thái, được điều phối bởi một tác tử LangGraph tự trị \cite{langgraph2024}: tầng này thực thi suy luận nhận thức chuyên sâu, truy xuất ngữ cảnh nền tảng từ công cụ Dual-RAG, đối sánh dữ liệu trên Bộ nhớ Đe dọa SQLite, và đưa ra các đề xuất ngăn chặn có cấu trúc như quy tắc chặn tường lửa động.
Thiết kế lai này giải quyết triệt để nút thắt độ trễ của LLM bằng cách giảm tải phần lớn lưu lượng lành tính ngay tại Tầng 1---tỷ lệ được định lượng thực nghiệm ở Chương~4---chỉ dành riêng tài nguyên suy luận LLM chuyên sâu cho các cuộc điều tra mối đe dọa phức tạp.

Quy trình hoạt động và luồng xử lý dữ liệu chi tiết theo thời gian thực của hệ thống được minh họa thông qua sơ đồ luồng chạy trong Hình~\ref{fig:sentinel_flow}.

\begin{figure}[!htbp]
    \centering
    \includegraphics[width=0.95\textwidth,height=0.82\textheight,keepaspectratio]{images/sentinel_flow.png}
    \caption{Sơ đồ luồng xử lý và luồng dữ liệu chi tiết theo thời gian thực của hệ thống SENTINEL. (Nguồn: Tác giả)}
    \label{fig:sentinel_flow}
\end{figure}

Như được trình bày trong sơ đồ luồng chạy, quy trình bắt đầu ngay khi nhật ký sự kiện mạng đầu vào được thu nạp. Bộ lọc Tầng 1 so khớp chữ ký qua Rule Engine; nếu khớp chữ ký WAF hoặc chèn mã, hệ thống đưa ra phán quyết \texttt{BLOCK\_IP} hoặc \texttt{ESCALATE}, còn lưu lượng trong whitelist thì bị loại bỏ. Trường hợp ngược lại, các thuộc tính số của log được cập nhật vào bộ tính toán Welford để xác định điểm Z-Score. Các dòng dữ liệu lành tính bị loại bỏ hoặc ghi vào nhật ký quyết định cục bộ, trong khi các dòng có điểm rủi ro vượt ngưỡng sẽ bị leo thang. Tầng 2 tiếp nhận cụm dữ liệu leo thang, trước tiên nén và đóng gói nonce (chạy thuật toán phân tích cú pháp log Drain \cite{he2017drain} để trích xuất khuôn mẫu sự kiện), sau đó \emph{luôn} truy xuất ngữ cảnh nền từ cơ sở tri thức Dual-RAG (được tăng tốc bằng Bộ đệm Ngữ nghĩa cho các khuôn mẫu lặp lại) cùng với Bộ nhớ Đe dọa dài hạn. Tác tử LangGraph sau đó kích hoạt nút phân loại LLM để đưa ra một phán quyết JSON có cấu trúc duy nhất. Sau khi ra quyết định, kết quả được kiểm định, làm sạch, ghi vào cơ sở dữ liệu với chuỗi HMAC an toàn và thực thi chặn tường lửa hoặc định tuyến cảnh báo tới giao diện HITL Streamlit.

\section{Định vị trong Ngăn xếp SOC Hiện có}
Một đánh giá công bằng phải phân tách rõ điều SENTINEL đóng góp với điều mà công cụ SOC trưởng thành đã cung cấp. Phát hiện dựa trên chữ ký và luật ở tốc độ đường truyền (Snort, Suricata, Zeek), tương quan khối lượng lớn (Splunk, QRadar, Elastic), và tự động hóa phản hồi tất định (các nền tảng SOAR như Cortex XSOAR, Splunk Phantom, hay Tines) đều đã là chuẩn mực từ lâu, và nhiều SIEM còn tự gán định danh MITRE ATT\&CK cho cảnh báo. SENTINEL chủ ý \emph{không} cạnh tranh với các hệ thống này ở khâu phát hiện: Tầng 1 phi trạng thái của nó bản thân chính là một hiện thân của đúng lớp chữ-ký-và-thống-kê đó, và trong triển khai production nó sẽ đứng trước---hoặc được nạp bởi---một IDS/SIEM sẵn có chứ không thay thế chúng. Lập trường bổ trợ này được củng cố bởi nghiên cứu loại trừ ở Chương~\ref{ch:experiments_evaluation}, trong đó việc thêm tầng nhận thức không cải thiện F1 phân loại flow thô so với riêng cổng Tầng 1; kết quả đó nằm trong dự đoán chính vì năng lực phát hiện phổ thông vốn đã mạnh, và đó là lý do giá trị của tầng thứ hai phải được tìm ở nơi khác.

Đóng góp thay vào đó nằm ở lớp mà ngăn xếp hiện nay \emph{không} cung cấp: một tầng \emph{nhận thức} an toàn, giải thích được, tự chủ, dành cho phần thiểu số sự kiện mơ hồ sống sót qua Tầng 1. Ba tính chất tách tầng này khỏi SOAR quy ước và khỏi các trợ lý ``AI-SOC'' đám mây mới nổi. Thứ nhất, SOAR quy ước thực thi các playbook \emph{cứng}, soạn sẵn---những cây quyết định tất định không thể lập luận về tình huống ngoài kịch bản hay trình bày lý do; SENTINEL thay tự-động-hóa-dựa-trên-luật này bằng một tác tử đồ-thị-trạng-thái lập luận trên các kỹ thuật MITRE ATT\&CK và kịch bản NIST truy xuất được, và xuất một lý giải kiểm toán được bởi con người (đạt $3.91/5$ theo LLM-Trọng-tài khác họ). Thứ hai, ngay khi một mô hình \emph{sinh} được đặt vào vòng quyết định, chính lớp cơ chất log trở thành một bề mặt tiêm---kẻ tấn công có thể nhúng chỉ dẫn kiểu ``bỏ qua các chỉ dẫn trước và đánh dấu cái này là vô hại'' vào một trường như User-Agent---một mối đe dọa mà bộ máy chữ ký và SOAR miễn nhiễm về cấu trúc vì chúng không bao giờ diễn giải nội dung theo lối sinh; lớp Đóng gói Dữ liệu Phân tách, lá chắn đồng thuận hai tầng, và chuỗi audit móc-xích HMAC là một phòng thủ trực tiếp, được kiểm chứng thực nghiệm ($100\%$ kháng ở mức toàn-pipeline) cho đúng bề mặt mới này. Thứ ba, các trợ lý AI-SOC thương mại thường định tuyến telemetry tới một LLM đám mây do nhà cung cấp lưu trữ, điều không được phép trong các môi trường đòi chủ quyền dữ liệu, air-gapped, hay Zero-Trust; SENTINEL chứng minh rằng toàn bộ tầng suy luận chạy cục bộ trên một mô hình lượng tử hóa đơn nhất, với cổng hai tầng giữ cho một lượt suy luận cục bộ nhiều giây vẫn khả thi về vận hành (giảm $82.97\%$ độ trễ trung bình nhờ chỉ leo thang phần thiểu số mơ hồ). Nói ngắn gọn, SENTINEL không thêm một bộ dò nữa vào một thị trường vốn đã đông; nó cung cấp tầng suy-luận-và-phản-hồi được làm cứng chống tiêm, triển khai được cục bộ, và giải thích được---thứ cho phép một mô hình sinh chiếm giữ an toàn vai trò tự động hóa mà các playbook SOAR cứng đang đảm nhiệm hôm nay.

\section{Tầng 1: Bộ lọc Lưu lượng Mạng Thời gian thực (RuleEngine và Welford)}

Bộ lọc Tầng 1 chịu trách nhiệm thu nạp và xử lý theo thời thực luồng nhật ký tốc độ cao. Để quản lý trạng thái kết nối mà không gây ra nút thắt cổ chai về đọc/ghi đĩa (disk I/O), bộ lọc tích hợp một trình quản lý đường cơ sở phiên tạm thời (temporal session baseline manager) để ánh xạ các địa chỉ IP nguồn đang hoạt động. Đối với mỗi phiên IP hoạt động, hệ thống tính toán các đặc trưng lưu lượng bằng thuật toán trực tuyến Welford \cite{welford1962}. Các công thức truy hồi cho giá trị trung bình động $\mu_k$ và tổng bình phương các độ lệch $M_{2,k}$ đã được trình bày tại Mục~\ref{sec:welford} (Phương trình~\ref{eq:welford_mean}--\ref{eq:welford_var}), cho ra phương sai cuộn $\sigma_k^2$ ổn định về số học trên một luồng liên tục các thuộc tính số của nhật ký (chẳng hạn như thời gian luồng, số lượng gói tin, hoặc kích thước byte). Đối với mỗi bản ghi nhật ký đầu vào, hệ thống sau đó tính toán điểm Z-Score $Z$ của các thuộc tính lưu lượng đặc trưng này:

\begin{equation}
    Z = \frac{x_k - \mu_k}{\sigma_k}
\end{equation}

Nếu điểm Z-Score vượt quá một ngưỡng tới hạn (được xác định trong cấu hình, ví dụ: $\alpha > 3.5\sigma$), gói tin sẽ được gắn nhãn hành động `ESCALATE` (leo thang) và chuyển tiếp lên Tầng 2. Để duy trì chuỗi sự kiện theo thời gian của các chiến dịch tấn công âm thầm và chậm rãi (low-and-slow), hệ thống tiếp tục cập nhật đường cơ sở của các phiên. Song song với các kiểm tra thống kê dựa trên thuật toán Welford, Tầng 1 vận hành một bộ máy luật tĩnh (\texttt{RuleEngine}) thực hiện so khớp mẫu chữ ký nhanh bằng biểu thức chính quy (regex) để ngăn chặn tức thì các mẫu câu lệnh đã biết (như cú pháp tấn công SQL Injection) và áp dụng các hành động tĩnh \texttt{DROP} (loại bỏ) hoặc \texttt{BLOCK\_IP} (chặn IP) mà không cần đi qua Tầng 2.

Về mặt cấu trúc phần mềm, Tầng 1 được định nghĩa thông qua hai lớp Python chính là \texttt{RunningStats} và \texttt{SessionBaseline}. Lớp \texttt{RunningStats} lưu trữ động số lượng, giá trị trung bình và tổng bình phương độ lệch trên bộ nhớ RAM, thực hiện các phép toán cập nhật thống kê trực tuyến với độ phức tạp thời gian tối ưu O(1). Để ngăn chặn tràn bộ nhớ khi lưu trữ trạng thái của vô số IP nguồn khác nhau, lớp \texttt{SessionBaseline} quản lý phiên thông qua cấu trúc bảng băm kết hợp với một luồng dọn rác (garbage collection daemon) định kỳ giải phóng các phiên không hoạt động dựa trên cấu hình TTL (Time-To-Live). Quy trình so khớp mẫu tĩnh được tăng tốc bằng các đối tượng biểu thức chính quy được biên dịch trước (pre-compiled regex) từ thư viện \texttt{re} của Python.

\section{Tầng 2: Tác tử AI Nhận thức LangGraph}

Các luồng bất thường được leo thang sẽ được tiếp nhận bởi công cụ nhận thức Tầng 2, được mô hình hóa dưới dạng Đồ thị có hướng không chu trình (Directed Acyclic Graph - DAG) do LangGraph \cite{langgraph2024} biên dịch và hoạt động như một Máy trạng thái hữu hạn (Finite State Machine - FSM). Dữ liệu phiên và ngữ cảnh suy luận được lưu giữ bên trong một đối tượng trạng thái trung tâm, \texttt{SentinelState}. Đồ thị đã biên dịch nối tiếp một nút rào chắn (nén và đóng gói phân tách), một nút ngữ cảnh RAG, một nút phân loại LLM, và một nút ánh xạ ATT\&CK có điều kiện theo trình tự, kết thúc tại một trong hai nút hành động (nút thực thi hành động hoặc nút chuyển con người phê duyệt). Nút ngữ cảnh RAG \emph{luôn} truy xuất ngữ cảnh nền từ cơ sở tri thức lai vector--từ khóa trước khi phân loại, chứ không phải chỉ được kích hoạt qua một nhánh điều kiện khi thiếu tri thức. Nút phân loại LLM sau đó đánh giá cụm nhật ký leo thang bằng mô hình ngôn ngữ lớn lượng tử hóa cục bộ Gemma 2 \cite{gemma_2_2024} và đưa ra một phán quyết có cấu trúc duy nhất. Mọi phán quyết mối-đe-doạ đáng-hành-động (\texttt{BLOCK\_IP}, \texttt{ALERT}, hoặc \texttt{AWAIT\_HITL}) được định tuyến qua nút ánh xạ ATT\&CK; nút này tất định neo vào kỹ thuật mà nút phân loại đã căn cứ và cấu trúc hóa nó thành một bản ghi MITRE hợp lệ theo lược đồ---tactic, technique, sub-technique (nếu có), URL chuẩn, điểm tin cậy ánh xạ, và phản hồi đề xuất suy ra từ tactic---trước khi phán quyết tiến tới nút hành động; các phán quyết lành tính \texttt{LOG} bỏ qua nút ánh xạ và kết thúc chu kỳ. Đồ thị định tuyến tất định dựa trên hành động đã được kiểm định này---phán quyết mức cao chuyển tới nút thực thi, phán quyết mơ hồ hoặc độ tin cậy thấp chuyển sang hàng đợi phê duyệt của con người, còn phán quyết lành tính thì kết thúc chu kỳ---mà không quay vòng lặp lại. Các trạng thái định tuyến và chuyển dịch được trình bày chi tiết trong Bảng~\ref{tab:agent_fsm_transitions}.

\begin{table}[!htbp]
    \centering
    \caption{Ma trận chuyển trạng thái của FSM Tác tử Nhận thức Tầng 2. (Nguồn: Tác giả)}
    \label{tab:agent_fsm_transitions}
    \begin{tabular}{|>{\raggedright\arraybackslash}p{0.23\textwidth}|>{\raggedright\arraybackslash}p{0.25\textwidth}|>{\raggedright\arraybackslash}p{0.19\textwidth}|>{\raggedright\arraybackslash}p{0.15\textwidth}|}
        \hline
        \textbf{Trạng thái hiện tại} & \textbf{Đầu vào / Điều kiện} & \textbf{Trạng thái tiếp theo} & \textbf{Hành động đầu ra} \\
        \hline
        \texttt{GUARDRAILS} & Nhận cụm log leo thang & \texttt{RAG\_CONTEXT} & Nén \& đóng gói nonce \\
        \hline
        \texttt{RAG\_CONTEXT} & Truy xuất ngữ cảnh thành công & \texttt{LLM\_TRIAGE} & Bổ sung ngữ cảnh MITRE\slash NIST \\
        \hline
        \texttt{LLM\_TRIAGE} & action $\in$ \{\texttt{BLOCK\_IP}, \texttt{ALERT}, \texttt{AWAIT\_HITL}\} & \texttt{ATTACK\_MAPPER} & Neo \& cấu trúc hóa bản ghi MITRE ATT\&CK \\
        \hline
        \texttt{LLM\_TRIAGE} & action $\in$ \{\texttt{LOG}, \texttt{DROP}\} & \texttt{END} & Kết thúc chu kỳ \\
        \hline
        \texttt{ATTACK\_MAPPER} & action $\in$ \{\texttt{BLOCK\_IP}, \texttt{ALERT}\} & \texttt{ACTION\_EXEC} & Thực thi chặn / phát cảnh báo \\
        \hline
        \texttt{ATTACK\_MAPPER} & action = \texttt{AWAIT\_HITL} (hoặc \texttt{BLOCK\_IP} với độ tin cậy $<0.5$) & \texttt{HITL} & Đưa vào hàng đợi chuyên viên SOC \\
        \hline
        \texttt{ACTION\_EXEC} / \texttt{HITL} & Hoàn tất hành động & \texttt{END} & Ghi bản ghi kiểm toán HMAC \\
        \hline
    \end{tabular}
\end{table}

Việc ràng buộc nút phân loại vào các tài liệu MITRE/NIST đã truy xuất giúp giới hạn LLM trong ngữ cảnh đã xác thực, giảm thiểu ảo giác (hallucination). Nút phân loại tạo ra một cấu trúc dữ liệu JSON chi tiết chứa phân loại mối đe dọa, điểm độ tin cậy và các quy tắc ngăn chặn được đề xuất. Các phán quyết độc hại có độ tin cậy cao được chuyển tới nút thực thi hành động, trong khi các phán quyết độ tin cậy thấp hoặc mơ hồ được định tuyến đến hàng đợi phê duyệt của con người (Human-in-the-Loop - HITL). Ngoài ra, một bộ \texttt{DecisionValidator} tất định sẽ hạ cấp mọi phán quyết \texttt{BLOCK\_IP} có độ tin cậy dưới 0.5 thành \texttt{AWAIT\_HITL}, và ghi đè LLM thành \texttt{AWAIT\_HITL} mỗi khi nó mâu thuẫn với tín hiệu tấn công của Tầng 1 (lá chắn đồng thuận Tầng 1/Tầng 2 chống thao túng ngữ nghĩa). Chuỗi kỹ thuật dạng văn bản tự do do nút phân loại sinh ra sau đó được nút ánh xạ ATT\&CK chuẩn hóa: với các lớp tấn công ứng dụng web kinh điển, một bảng tra cứu được người soạn kiểm chứng sẽ giải mã kỹ thuật một cách tất định; còn các trường hợp khác, nút neo vào định danh kỹ thuật hợp lệ mà nút phân loại đã sinh và chỉ lui về truy xuất RRF khi định danh đó vắng mặt---nên nút ánh xạ cấu trúc hóa chứ không ghi đè phán quyết của bộ phân loại.

Hiện thực hóa kỹ thuật của đồ thị trạng thái sử dụng trình biên dịch của LangGraph. Trạng thái của tác tử được định nghĩa như một lớp dữ liệu có cấu trúc (dataclass \texttt{SentinelState}), trong đó các trường lưu trữ khóa như chỉ số IOC và lịch sử quyết định được thiết lập ở chế độ chỉ cho phép ghi thêm (append-only) nhằm chống lại sự trôi dạt ngữ nghĩa (semantic drift) trong quá trình suy luận. Việc kết nối LLM được quản lý bởi lớp \texttt{LLMClient} giao tiếp với máy chủ llama.cpp chạy mô hình Gemma 2 \cite{gemma_2_2024} cục bộ, cấu hình với các tham số khắt khe như \texttt{temperature=0.1} và giới hạn cửa sổ ngữ cảnh cố định 8192 token để bảo đảm suy luận gần như tất định và ngăn chặn quá tải bộ nhớ VRAM.

\section{Cơ chế Sinh Tăng cường Truy xuất Kép (Dual-RAG) và Bộ nhớ Ngữ nghĩa}

Cơ chế Dual-RAG định hình các quyết định của tác tử bằng cách truy xuất ngữ cảnh từ hai cơ sở dữ liệu chuyên ngành thống nhất: cơ sở tri thức MITRE ATT\&CK \cite{strom2018mitre} dùng để ánh xạ các hành vi, chiến thuật và kỹ thuật đe dọa (TTPs), và kho lưu trữ kịch bản NIST SP 800-61r2 cung cấp các quy trình xử lý sự cố tiêu chuẩn \cite{nist80061}. Thay vì chỉ dựa vào khớp từ khóa đơn thuần hoặc các vector ngữ nghĩa dày đặc một cách độc lập \cite{rag2020}, quá trình truy xuất sử dụng một chiến lược tìm kiếm lai (hybrid search). Bảng so sánh giữa các thành phần tìm kiếm dày đặc (dense) và thưa thớt (sparse) được chi tiết hóa trong Bảng~\ref{tab:search_comparison}.

\begin{table}[!htbp]
    \centering
    \caption{So sánh các mô hình truy xuất trong chiến lược tìm kiếm lai. (Nguồn: Tác giả)}
    \label{tab:search_comparison}
    \begin{tabular}{|>{\raggedright\arraybackslash}p{0.22\textwidth}|>{\raggedright\arraybackslash}p{0.30\textwidth}|>{\raggedright\arraybackslash}p{0.30\textwidth}|}
        \hline
        \textbf{Khía cạnh} & \textbf{Tìm kiếm Vector Dày đặc (FAISS)} & \textbf{Tìm kiếm Từ khóa Thưa thớt (BM25)} \\
        \hline
        \textbf{Mô hình nền tảng} & all-MiniLM-L6-v2, embedding 384 chiều \cite{reimers2019sentence} & Trọng số từ khóa BM25 \\
        \hline
        \textbf{Logic so khớp} & Độ tương đồng Cosine ngữ nghĩa \cite{faiss2019} & Tần suất từ - Tần suất nghịch đảo tài liệu (TF-IDF) \cite{bm25_2009} \\
        \hline
        \textbf{Ưu điểm} & Bắt được các từ đồng nghĩa, khái niệm và ngữ nghĩa trừu tượng & So khớp chính xác các định danh kỹ thuật (CVEs, IPs, số cổng) \\
        \hline
        \textbf{Hạn chế} & Yếu khi khớp chính xác các từ khóa băm kỹ thuật & Bị mù trước các từ đồng nghĩa ngữ nghĩa và ngữ cảnh khái niệm \\
        \hline
    \end{tabular}
\end{table}

Điểm số từ các luồng xử lý dày đặc (FAISS) và thưa thớt (BM25) được tổng hợp lại bằng thuật toán Tổng hợp xếp hạng nghịch đảo (Reciprocal Rank Fusion - RRF) \cite{rrf2009}, tính toán một điểm số thống nhất $R$ cho mỗi tài liệu $d$ trên tập hợp các lượt tìm kiếm $M$:

\begin{equation}
    R(d) = \sum_{m \in M} \frac{1}{k + r_m(d)}
\end{equation}

trong đó $r_m(d)$ là thứ hạng của tài liệu $d$ trong lượt tìm kiếm $m$, và $k$ là một hằng số tỷ lệ (thường chọn $k=60$). Ba đoạn văn bản có thứ hạng cao nhất từ danh sách hợp nhất này được chuyển tới nút phân loại làm ngữ cảnh nền tảng. Để tránh chi phí nhúng và tìm kiếm dư thừa, một mô-đun Bộ đệm Ngữ nghĩa (Semantic Cache) được triển khai phía trước luồng truy xuất. Bộ đệm này sinh khóa tất định bằng cách tính mã băm SHA-256 của khuôn mẫu nhật ký đã khai thác, nhờ đó các log leo thang có cùng khuôn mẫu tấn công sẽ tái sử dụng kết quả truy xuất đã tính trước thay vì nhúng và tìm kiếm lại, rút ngắn thời gian xử lý mỗi truy vấn từ hàng trăm mili giây xuống còn micro giây. Bộ đệm được hiện thực dưới dạng cấu trúc LRU (Least Recently Used) trong bộ nhớ (một \texttt{OrderedDict} của Python) giới hạn 500 mục, kèm thời gian sống (TTL) 1800 giây để loại bỏ các khuôn mẫu cũ.

\section{Bộ nhớ Đe dọa và Liên kết Chiến dịch APT}

Các tác nhân đe dọa tinh vi thường chia nhỏ chuỗi tấn công của chúng thành các sự kiện có mức độ nghiêm trọng thấp và phân tán trên các khoảng thời gian dài nhằm vượt qua các bộ lọc gói tin phi trạng thái. Để vạch trần các chiến dịch tấn công âm thầm và chậm rãi này \cite{apt_stealth2021}, kiến trúc SENTINEL tích hợp một mô-đun Bộ nhớ Đe dọa (Threat Memory) bền vững, được triển khai dưới dạng cơ sở dữ liệu SQLite cục bộ. Mô-đun này được xây dựng trên các lược đồ (schema) chuyên dụng để theo dõi các cảnh báo lịch sử, điểm uy tín của các thực thể, và chuỗi tấn công theo thời gian. Chi tiết về các lược đồ cơ sở dữ liệu được trình bày trong Bảng~\ref{tab:threat_memory_schema}.

\begin{table}[!htbp]
    \centering
    \caption{Thông số kỹ thuật lược đồ cơ sở dữ liệu Bộ nhớ Đe dọa. (Nguồn: Tác giả)}
    \label{tab:threat_memory_schema}
    \begin{tabular}{|>{\raggedright\arraybackslash}p{0.25\textwidth}|>{\raggedright\arraybackslash}p{0.33\textwidth}|>{\raggedright\arraybackslash}p{0.28\textwidth}|}
        \hline
        \textbf{Tên bảng} & \textbf{Trường khóa chính / thuộc tính} & \textbf{Mục đích chức năng} \\
        \hline
        \texttt{ip\_reputation} & \texttt{ip} (PK), \texttt{reputation\_score}, \texttt{total\_incidents}, \texttt{total\_blocks}, \texttt{first\_seen}, \texttt{last\_seen}, \texttt{last\_mitre\_technique} & Theo dõi điểm uy tín dài hạn (0--100) và lịch sử sự cố của các IP nguồn, có cơ chế suy giảm theo thời gian. \\
        \hline
        \texttt{known\_entities} & \texttt{id} (PK), \texttt{entity\_type}, \texttt{entity\_value} (duy nhất), \texttt{description}, \texttt{is\_active} & Khởi tạo trước các công cụ/dịch vụ nội bộ đã xác minh (như AD, máy quét) để giảm thiểu cảnh báo giả. \\
        \hline
        \texttt{threat\_events} & \texttt{id} (PK), \texttt{src\_ip}, \texttt{dst\_ip}, \texttt{apt\_phase}, \texttt{apt\_day}, \texttt{label}, \texttt{timestamp} & Ghi nhận từng sự kiện leo thang để liên kết chuỗi APT đa ngày nổi lên dần. \\
        \hline
        \texttt{apt\_indicators} & \texttt{id} (PK), \texttt{indicator\_type}, \texttt{indicator\_value}, \texttt{confidence}, \texttt{occurrence\_count}, \texttt{mitre\_chain} & Lưu các chỉ dấu liên kết APT đã xác nhận phục vụ giám sát chiến dịch dài hạn. \\
        \hline
    \end{tabular}
\end{table}

Điểm uy tín của một địa chỉ IP sẽ tăng lên dựa trên mức độ nghiêm trọng của các cảnh báo mà nó kích hoạt, giới hạn tối đa là 100. Ngược lại, một hàm suy giảm (decay function) định kỳ chạy để giảm dần điểm uy tín của các IP không còn hoạt động:

\begin{equation}
    S_{new} = S_{old} \times (1 - \lambda)^{t}
\end{equation}

trong đó $\lambda$ là tỷ lệ suy giảm và $t$ là thời gian không hoạt động đã trôi qua. Quy trình liên kết APT sẽ quét bảng \texttt{threat\_events}. Nếu phát hiện một IP nguồn kích hoạt các dị thường trong ít nhất hai ngày khác nhau (có thể trải trên nhiều giai đoạn của chuỗi kill-chain), hệ thống sẽ tự động nâng mức phân loại mối đe dọa lên nghiêm trọng cao, xem đây là một chiến dịch đa ngày nổi lên dần thay vì một sự kiện đơn lẻ. Trong hiện thực, các khoảng cách thời gian và điểm tích lũy uy tín được tính bằng các truy vấn upsert và tổng hợp phía Python trên kho SQLite, tránh việc phải lưu giữ trạng thái dài hạn trong bộ nhớ làm việc của tác tử.

\section{Tầng Rào chắn Bảo mật AI (Guardrails)}

Do các LLM dễ bị tấn công chiếm quyền điều khiển chỉ thị (instruction hijacking / prompt injection), hệ thống SENTINEL triển khai một tầng Rào chắn Bảo mật AI (AI Security Guardrails) nhiều lớp nhằm cách ly công cụ suy luận khỏi các dữ liệu đầu vào không đáng tin cậy \cite{promptinjection2023}. Cơ chế phòng thủ cốt lõi dựa trên kỹ thuật Đóng gói Dữ liệu Phân tách (Delimited Data Encapsulation). Đối với mỗi phiên làm việc, hệ thống tự động tạo ra một mã token thập lục phân ngẫu nhiên bảo mật bằng mật mã (ví dụ: \texttt{<|DATA\_8F3E|>}) bằng hàm \texttt{secrets.token\_hex()} của Python (dựa trên \texttt{os.urandom}). Phần dữ liệu nhật ký thô chưa được kiểm chứng sẽ được bao bọc bên trong các thẻ phân tách này trong khuôn mẫu prompt, chỉ thị cho LLM rằng bất kỳ mã lệnh nào xuất hiện bên trong ranh giới này đều phải được xử lý nghiêm ngặt dưới dạng dữ liệu tĩnh thụ động thay vì các chỉ thị thực thi.

Để chống lại các cuộc tấn công Từ chối dịch vụ Token (Token DoS) và tràn ngữ cảnh, hệ thống tích hợp một trình Khai thác Mẫu Nhật ký (Log Template Miner) dựa trên thuật toán phân tích cú pháp log Drain \cite{he2017drain}. Trình khai thác này nhóm các bản ghi nhật ký tương đồng thành các khuôn mẫu chung, loại bỏ các tham số ngẫu nhiên (chẳng hạn như mốc thời gian động hoặc ID luồng). Hoạt động trong ngân sách 4000 token, Bộ Quản lý Ngân sách Token (Token Budget Manager) sau đó phân bổ tối đa 60\% ngân sách này cho các khuôn mẫu đã nén và 40\% còn lại cho các điểm dị biệt có entropy cao. Cuối cùng, một Bộ lọc Làm sạch Đầu ra (Output Sanitization Filter) áp dụng cơ chế xác thực lược đồ JSON nghiêm ngặt, kiểm chứng câu trả lời của LLM chỉ chứa các trường được phép, đồng thời loại bỏ các thẻ hình ảnh markdown, kịch bản SVG, và các vector vượt rào mã hóa thập lục phân \cite{watchtower2026}. Lớp này được xây dựng dạng module độc lập \texttt{output\_sanitizer} để lọc dữ liệu trước khi nạp vào SQLite hoặc giao diện dashboard.

\section{Toàn vẹn Dữ liệu và Bảo vệ bằng Chuỗi HMAC}

Để bảo vệ cơ sở dữ liệu nhật ký bảo mật trước các hành vi giả mạo bởi người nội bộ độc hại hoặc những kẻ tấn công giành được quyền truy cập vào hệ thống tệp tin, hệ thống triển khai kỹ thuật chuỗi liên kết log bằng HMAC (HMAC log chaining)~\cite{hmac1997}. Mỗi bản ghi nhật ký an ninh được liên kết mật mã với bản ghi liền trước nó, tạo thành một chuỗi vết kiểm toán tương tự cấu trúc blockchain nội bộ. Đối với bản ghi thứ $i$, mã băm bản ghi $H_i$ được tính toán như sau:

\begin{equation}
    H_i = \text{HMAC-SHA256}(D_i \parallel H_{i-1}, K)
\end{equation}

trong đó $D_i$ là dữ liệu tải tin của bản ghi hiện tại, $H_{i-1}$ là mã băm của bản ghi liền trước, và $K$ là khóa bí mật HMAC được lưu trữ an toàn. Nếu kẻ tấn công cố tình thay đổi nhãn nhật ký hoặc xóa một bản ghi, chuỗi băm sẽ bị đứt gãy và các quy trình kiểm tra xác thực chạy trên bảng điều khiển sẽ lập tức gắn cờ phát hiện hành vi giả mạo. Ngoài ra, hệ thống cũng bảo vệ cơ sở dữ liệu vector RAG trước các cuộc tấn công đầu độc ngữ nghĩa (semantic poisoning). Các tệp tin cơ sở dữ liệu vector được giám sát chặt chẽ bởi một trình quản lý mã checksum chuyên dụng thực hiện tính toán và so sánh giá trị băm SHA-256 của các tệp tin trước khi nạp chúng vào bộ nhớ, đảm bảo không có tài liệu độc hại nào bị tiêm vào bộ nhớ truy xuất của tác tử.

Việc liên kết chuỗi toàn vẹn được hiện thực trong quy trình ghi nhật ký kiểm toán, còn lớp \texttt{ActionValidator} thực thi danh sách cho phép (allowlist) các hành động an ninh hợp lệ, phát hiện và từ chối mọi siêu ký tự shell trong các trường của tải tin đầu vào nhằm ngăn chặn nguy cơ thực thi mã từ xa (RCE) trong quá trình điều phối kịch bản cấu hình tường lửa.

\section{Giao diện Streamlit Bảng điều khiển Giám sát SOC (Streamlit UI)}

Giao diện Bảng điều khiển Giám sát SOC được phát triển trên Streamlit, tuân thủ mô hình phản hồi (reactive model) liên lạc liên tục với backend và cơ sở dữ liệu SQLite. Streamlit UI tích hợp ba phân hệ hiển thị chính. Thứ nhất, phân hệ luồng và hàng đợi hiển thị trực quan thông lượng log đi qua Tầng 1 và các kết quả phân tích phân loại của Tầng 2. Thứ hai, phân hệ phân tích bộ nhớ đe dọa vẽ sơ đồ chuỗi lan truyền cảnh báo APT của các địa chỉ IP theo dòng thời gian. Thứ ba, phân hệ kiểm tra toàn vẹn log liên tục xác thực chuỗi băm HMAC của cơ sở dữ liệu nhật ký và đưa ra các đèn tín hiệu cảnh báo (xanh: toàn vẹn; đỏ: đã bị can thiệp/sửa đổi). Ngoài ra, giao diện hỗ trợ cơ chế phê duyệt HITL, cho phép phân tích viên kiểm tra các đề xuất của tác tử và nhấn nút đồng ý/hủy bỏ cập nhật cấu hình tường lửa thực tế.

\section{Kiến trúc ảo hóa Docker và Triển khai hạ tầng}

Hệ thống được đóng gói thông qua Docker với cấu hình \texttt{Dockerfile} đa tầng (multi-stage build) để tối ưu dung lượng ảnh sản phẩm. Trình biên dịch llama.cpp trong container được kích hoạt CUDA để tận dụng GPU của hệ thống vật lý giúp đẩy nhanh tốc độ suy luận của LLM. Tệp điều phối \texttt{docker-compose.yml} thiết lập mạng nội bộ liên kết container tác tử và container Streamlit, chia sẻ chung các phân vùng ổ đĩa (volume mounts) chứa các cơ sở dữ liệu SQLite là \texttt{threat\_memory.db} và \texttt{guardrails\_audit.db} để đảm bảo tính lưu trữ bền vững độc lập với vòng đời của các container.

\section{Khả năng Mở rộng và Suy biến Nhẹ nhàng}

Một phản biện tự nhiên đối với mọi đường ống có LLM trong vòng lặp là về thông lượng: nếu dữ liệu đo lường đổ về hàng tỷ sự kiện mỗi ngày, liệu mô hình ngôn ngữ có trở thành nút thắt cổ chai không thể vượt qua? Thiết kế hai tầng trả lời câu hỏi này bằng cách coi hai tầng là hai chế độ mở rộng khác nhau về bản chất, thay vì trang bị phần cứng suy luận ngày càng lớn hơn.

\textbf{Tầng 1 là bài toán thông lượng, không phải bức tường độ trễ.} Vì mỗi lần cập nhật Welford và so khớp regex chạy trong thời gian $O(1)$ và trạng thái phiên theo từng IP là độc lập lẫn nhau, Tầng 1 song song hóa gần như hoàn hảo: khối lượng công việc được phân mảnh sạch sẽ trên các lõi và máy chủ theo hàm băm của IP nguồn. Trong prototype Python, mỗi quyết định mất khoảng $0.6$\,ms, nên một lõi đơn duy trì cỡ một nghìn sự kiện mỗi giây; thông lượng một tỷ sự kiện mỗi ngày trung bình khoảng $1.2\times10^{4}$ sự kiện mỗi giây, nằm trong tầm với của một nút nhiều lõi vừa phải, với các đợt bùng nổ được hấp thụ bằng phân mảnh theo chiều ngang. Một bản hiện thực ở dạng biên dịch hoặc xử lý luồng (ví dụ một probe eBPF hoặc một toán tử luồng viết bằng Rust) nâng trần thông lượng mỗi lõi lên nhiều bậc độ lớn. Do đó, chính thuật toán $O(1)$ mới là thứ giúp đạt được tốc độ đường truyền; prototype kiểm chứng tỷ lệ lọc và tính đúng đắn chứ không phải thông lượng line-rate.

\textbf{Tầng nhận thức không bao giờ thấy số lượng leo thang thô.} Mặc dù Tầng 1 chỉ chuyển tiếp phần thiểu số bất thường của lưu lượng vốn áp đảo là lành tính, ba lớp tiếp theo còn ép con số này nhỏ lại trước khi bất kỳ suy luận nào diễn ra. Thứ nhất, trình khai phá khuôn mẫu Drain nén các dòng log gần như giống hệt nhau---phần lặp lại chiếm đa số của các tấn công thực như quét cổng, dò mật khẩu, và lũ thể tích---thành một tập nhỏ các khuôn mẫu. Thứ hai, Bộ đệm Ngữ nghĩa, lập khóa bằng mã băm SHA-256 của mỗi khuôn mẫu, trả về phán quyết đã tính trước cho bất kỳ khuôn mẫu lặp lại nào mà không gọi lại mô hình. Thứ ba, một khi một IP bị chặn, bản ghi uy tín trong Bộ nhớ Đe dọa khiến các sự kiện tiếp theo của nó bị loại bỏ ngay tại Tầng 1 thay vì leo thang lại. Vì vậy, tải suy luận thực tế bị chặn trên không phải bởi khối lượng sự kiện thô mà bởi số khuôn mẫu bất thường \emph{thực sự mới}, vốn nhỏ hơn rất nhiều và được xử lý theo lô.

\textbf{Đơn vị suy luận là cả lô, không phải từng sự kiện.} Các log leo thang được gom vào một lô của mỗi chu kỳ và đưa vào mô hình trong \emph{một} lần gọi suy luận, thay vì mỗi sự kiện một lần gọi; số log mang theo mỗi lô chỉ bị giới hạn bởi ngân sách token, nên hàng chục bất thường đã nén được phân xử bởi một lần suy luận duy nhất ($\approx$5.7\,s). Quan trọng không kém, tầng nhận thức nằm \emph{ngoài} đường chặn đồng bộ: Tầng 1 đã đưa ra phán quyết tất định \texttt{BLOCK\_IP} hoặc \texttt{DROP} cho các tấn công rõ ràng ở tốc độ đường truyền, nên lô LLM chỉ đóng góp phần giải thích, nền tảng MITRE/NIST và tương quan chéo sự kiện một cách \emph{bất đồng bộ}. Do đó, một backlog ở tầng nhận thức chỉ làm chậm việc làm giàu chứ không làm chậm việc bảo vệ---việc giảm thiểu các mối đe dọa rõ ràng không bao giờ phụ thuộc vào độ trễ của LLM. Thông lượng suy luận sâu bền vững của một bộ tăng tốc vì vậy bằng (số log mỗi lô) chia cho (độ trễ mỗi lô), vượt xa tốc độ ngây thơ theo từng sự kiện đúng bằng hệ số lô, chưa kể các lần trúng cache.

\textbf{Hành vi dưới một trận lũ đối kháng.} Trường hợp tệ nhất đối với tầng nhận thức là một trận lũ các bất thường đều mới toàn bộ, khiến việc nén khuôn mẫu và đệm cache không thể hấp thụ. Hai nguyên lý thiết kế giới hạn thiệt hại. Thứ nhất, bản thân trận lũ đó là một tín hiệu mà Tầng 1 phát hiện được: sự tăng vọt tốc độ sự kiện là một điểm dị biệt Welford, kích hoạt một biện pháp giảm thiểu thô theo thể tích---ví dụ giới hạn tốc độ hoặc chặn các dải địa chỉ vi phạm---mà không cần suy luận từng sự kiện. Thứ hai, tầng nhận thức là một tài nguyên có dung lượng hữu hạn buộc phải đổ tải một cách nhẹ nhàng: khi hàng đợi của nó bão hòa, đường ống tụt về phán quyết tất định của Tầng 1 (chặn hoặc cảnh báo từ Z-score và bộ máy luật), chỉ hoãn lại phần giải thích và nền tảng phong phú hơn của LLM. Nhờ đó phạm vi bảo vệ an ninh được giữ vững dưới quá tải, và chỉ phần làm giàu nhận thức bị suy biến. Hệ quả là dung lượng GPU được trang bị cho dòng rỉ rả ổn định của các bất thường mới---vốn một bộ tăng tốc cấp phổ thông đã phục vụ đủ trong công trình này---thay vì cho line-rate tệ nhất, bởi việc scale mô hình cho bằng tốc độ thu nạp thô vừa tự phủ định kiến trúc vừa phi lý về kinh tế.

Cần nêu một cách trung thực rằng prototype hiện tại đã hiện thực tỷ lệ đổ tải và các lớp đệm cache nhưng chưa có một bộ điều khiển backpressure-và-suy-biến tường minh; dưới một trận lũ bất thường-mới kéo dài, nó sẽ làm bão hòa hàng đợi leo thang, như đã thừa nhận trong phần giới hạn ở Chương~\ref{ch:conclusion}. Việc hình thức hóa chính sách suy biến nhẹ nhàng này, cùng với việc scale GPU theo chiều ngang và phân rã đa tác tử nêu trong hướng phát triển tương lai, là con đường tự nhiên hướng tới một triển khai cấp sản xuất.

\section{Tóm tắt chương}

Chương 3 đã xây dựng toàn diện thiết kế kiến trúc và quy trình triển khai kỹ thuật của hệ thống SENTINEL. Bằng cách gộp bộ lọc mạng tốc độ đường truyền phi trạng thái (Tầng 1) với tác tử suy luận nhận thức có trạng thái (Tầng 2), đồng thời bổ sung các rào chắn AI cùng cơ chế chữ ký HMAC, hệ thống đạt được hiệu năng cao và độ bảo mật tuyệt đối trước các tấn công đối kháng. Quy trình tự động hóa được kiểm soát an toàn bởi con người qua hàng đợi HITL trên Streamlit UI và được đóng gói hoàn chỉnh trong Docker. Chương 4 tiếp theo sẽ mô tả chi tiết các thiết lập thực nghiệm, nghiên cứu loại trừ (ablation study) và các báo cáo kiểm định thống kê đánh giá hiệu năng thực tế của hệ thống.
